% ============================================================
\chapter{Transform\'ee de Fourier et \'EDP}
\label{ch:transformee_fourier}
% ============================================================

En 1807, Joseph Fourier affirme que toute fonction peut s'\'ecrire comme somme de sinuso\"ides. Lagrange est sceptique, Poisson dubitatif, mais l'id\'ee est trop puissante pour \^etre ignor\'ee. La \emph{transform\'ee de Fourier}, g\'en\'eralisation continue des s\'eries de Fourier, convertit une fonction de la variable spatiale en une fonction de la fr\'equence. Pour les EDP \`a coefficients constants, cette transformation est miraculeuse~: les d\'eriv\'ees deviennent des multiplications, les \'equations aux d\'eriv\'ees partielles se transforment en \'equations alg\'ebriques. L'\'equation de la chaleur, l'\'equation des ondes, l'\'equation de Schr\"odinger --- toutes se r\'esolvent \'el\'egamment par Fourier sur $\R^n$. Ce chapitre d\'eveloppe la transform\'ee de Fourier sur $L^1$ et $L^2$, \'etablit la formule d'inversion et le th\'eor\`eme de Plancherel, puis les applique syst\'ematiquement aux EDP classiques.

\section{Transform\'ee de Fourier sur $L^1(\R^n)$}

\begin{definition}[Transform\'ee de Fourier]
\index{transform\'ee de Fourier}
Soit $f \in L^1(\R^n)$. La \emph{transform\'ee de Fourier} de $f$ est la fonction
$\hat{f} : \R^n \to \C$ d\'efinie par~:
\[
    \hat{f}(\xi) = \mathcal{F}[f](\xi) = \int_{\R^n} f(x)\, e^{-2\pi i \inner{x}{\xi}} \, dx.
\]
\end{definition}

\begin{remarque}
Il existe diff\'erentes conventions pour la transform\'ee de Fourier. On utilise ici
la convention sym\'etrique avec $2\pi$ dans l'exponentielle, qui \'elimine les facteurs
$(2\pi)^{n/2}$ de la formule d'inversion.
\end{remarque}

\begin{proposition}[Propri\'et\'es de base]
\label{prop:proprietes_fourier}
Soient $f, g \in L^1(\R^n)$~:
\begin{enumerate}
    \item \textbf{Lin\'earit\'e}~: $\widehat{\alpha f + \beta g} = \alpha \hat{f} + \beta \hat{g}$.
    \item \textbf{Translation}~: $\widehat{f(\cdot - a)}(\xi) = e^{-2\pi i \inner{a}{\xi}} \hat{f}(\xi)$.
    \item \textbf{Modulation}~: $\widehat{e^{2\pi i \inner{a}{\cdot}} f}(\xi) = \hat{f}(\xi - a)$.
    \item \textbf{Dilatation}~: $\widehat{f(\lambda \cdot)}(\xi) = \frac{1}{\abs{\lambda}^n} \hat{f}(\xi/\lambda)$ pour $\lambda \neq 0$.
    \item \textbf{Convolution}~: $\widehat{f * g} = \hat{f} \cdot \hat{g}$.
\end{enumerate}
\end{proposition}

\begin{theoreme}[Lemme de Riemann-Lebesgue]
\index{lemme de Riemann-Lebesgue}
Si $f \in L^1(\R^n)$, alors $\hat{f} \in C_0(\R^n)$ (continue et tendant vers $0$
\`a l'infini)~: $\hat{f}(\xi) \to 0$ quand $\abs{\xi} \to \infty$.
\end{theoreme}

\begin{proof}
Pour $f$ indicatrice d'un rectangle, le r\'esultat est un calcul direct. Par
densit\'e des fonctions en escalier dans $L^1$ et par continuit\'e de $\mathcal{F}$
(car $\abs{\hat{f}(\xi)} \leq \norm{f}_{L^1}$), on \'etend \`a tout $f \in L^1$.
\end{proof}

\begin{theoreme}[Transform\'ee de la d\'eriv\'ee]
\index{transform\'ee de Fourier!d\'eriv\'ee}
Si $f$ et $\partial^\alpha f$ sont dans $L^1(\R^n)$, alors~:
\[
    \widehat{\partial^\alpha f}(\xi) = (2\pi i \xi)^\alpha \hat{f}(\xi),
\]
o\`u $\alpha$ est un multi-indice et $(2\pi i \xi)^\alpha = \prod_{j=1}^{n}(2\pi i \xi_j)^{\alpha_j}$.
\end{theoreme}

\begin{intuition}
La transform\'ee de Fourier convertit les d\'eriv\'ees en multiplications par des
polyn\^omes. C'est la raison fondamentale de son efficacit\'e pour r\'esoudre les
EDP lin\'eaires \`a coefficients constants~: une EDP devient une \'equation
alg\'ebrique dans l'espace de Fourier.
\end{intuition}

\section{Transform\'ee de Fourier sur $L^2(\R^n)$}

\begin{theoreme}[Plancherel]
\label{thm:plancherel}
\index{th\'eor\`eme de Plancherel}
La transform\'ee de Fourier s'\'etend en un isomorphisme isom\'etrique de $L^2(\R^n)$~:
\[
    \norm{\hat{f}}_{L^2} = \norm{f}_{L^2} \quad \forall\, f \in L^2(\R^n).
\]
Plus g\'en\'eralement, $\inner{\hat{f}}{\hat{g}}_{L^2} = \inner{f}{g}_{L^2}$
(identit\'e de Parseval).
\end{theoreme}

\begin{proof}[Id\'ee de preuve]
On montre d'abord le r\'esultat pour $f \in L^1 \cap L^2$ (espace dense dans $L^2$),
en calculant $\int \abs{\hat{f}}^2$ via la convolution $f * \bar{f}(-\cdot)$.
On \'etend ensuite par densit\'e et continuit\'e.
\end{proof}

\begin{theoreme}[Formule d'inversion]
\index{formule d'inversion de Fourier}
Si $f, \hat{f} \in L^1(\R^n)$, alors~:
\[
    f(x) = \int_{\R^n} \hat{f}(\xi)\, e^{2\pi i \inner{x}{\xi}} \, d\xi
    \quad \text{p.p.}
\]
\end{theoreme}

\section{Distributions temp\'er\'ees}

\begin{definition}[Espace de Schwartz]
\index{espace de Schwartz}
L'\emph{espace de Schwartz} $\mathcal{S}(\R^n)$ est l'espace des fonctions
$\varphi \in C^\infty(\R^n)$ dont toutes les d\'eriv\'ees d\'ecroissent plus vite
que tout polyn\^ome~:
\[
    \sup_{x \in \R^n} \abs{x^\alpha \partial^\beta \varphi(x)} < \infty
    \quad \forall\, \alpha, \beta \in \N^n.
\]
\end{definition}

\begin{definition}[Distribution temp\'er\'ee]
\index{distribution temp\'er\'ee}
L'espace $\mathcal{S}'(\R^n)$ des \emph{distributions temp\'er\'ees} est le dual
topologique de $\mathcal{S}(\R^n)$. La transform\'ee de Fourier s'\'etend \`a
$\mathcal{S}'$ par dualit\'e~:
\[
    \langle \hat{T}, \varphi \rangle = \langle T, \hat{\varphi} \rangle
    \quad \forall\, T \in \mathcal{S}', \; \varphi \in \mathcal{S}.
\]
\end{definition}

\begin{exemple}[Transform\'ee de la masse de Dirac]
$\hat{\delta} = 1$, car $\langle \hat{\delta}, \varphi \rangle = \langle \delta, \hat{\varphi} \rangle
= \hat{\varphi}(0) = \int \varphi(\xi)\, d\xi = \langle 1, \varphi \rangle$.
\end{exemple}

\begin{exemple}[Transform\'ee de la gaussienne]
\index{gaussienne}
Pour $f(x) = e^{-\pi \abs{x}^2}$, on a $\hat{f}(\xi) = e^{-\pi \abs{\xi}^2}$.
La gaussienne est un point fixe de la transform\'ee de Fourier.
\end{exemple}

\section{R\'esolution de l'\'equation de la chaleur}

Consid\'erons le probl\`eme de Cauchy pour l'\'equation de la chaleur~:
\begin{equation}\label{eq:chaleur_cauchy}
    u_t = \kappa\, \Delta u, \quad x \in \R^n,\; t > 0, \qquad
    u(x, 0) = u_0(x).
\end{equation}

\begin{theoreme}[Solution par transform\'ee de Fourier]
\label{thm:chaleur_fourier}
La solution de \eqref{eq:chaleur_cauchy} est donn\'ee par~:
\[
    u(x, t) = (G_t * u_0)(x) = \int_{\R^n} G_t(x - y)\, u_0(y)\, dy,
\]
o\`u $G_t$ est le \emph{noyau de la chaleur}~:
\[
    G_t(x) = \frac{1}{(4\pi \kappa t)^{n/2}} \exp\!\left(-\frac{\abs{x}^2}{4\kappa t}\right).
\]
\end{theoreme}

\begin{proof}
En appliquant la transform\'ee de Fourier en $x$~:
\[
    \hat{u}_t(\xi, t) = -4\pi^2 \kappa \abs{\xi}^2 \hat{u}(\xi, t).
\]
C'est une EDO en $t$, de solution $\hat{u}(\xi, t) = \hat{u}_0(\xi)\, e^{-4\pi^2 \kappa \abs{\xi}^2 t}$.
On reconna\^it $e^{-4\pi^2 \kappa \abs{\xi}^2 t} = \hat{G}_t(\xi)$, donc
$u = G_t * u_0$ par la propri\'et\'e de convolution.
\end{proof}

\begin{attention}[title=\textbf{Propagation instantan\'ee}]
Le noyau de la chaleur $G_t(x) > 0$ pour tout $x$ et tout $t > 0$. Ainsi,
m\^eme si $u_0$ est \`a support compact, $u(x, t) > 0$ pour tout $x$ d\`es que
$t > 0$~: la chaleur se propage \`a vitesse infinie.
\end{attention}

\section{R\'esolution de l'\'equation des ondes}

\begin{theoreme}[Solution de l'\'equation des ondes en dimension $n$]
Le probl\`eme de Cauchy~:
\[
    u_{tt} = c^2 \Delta u, \quad u(x,0) = u_0(x), \quad u_t(x,0) = v_0(x),
\]
a pour solution (dans l'espace de Fourier)~:
\[
    \hat{u}(\xi, t) = \hat{u}_0(\xi) \cos(2\pi c \abs{\xi} t)
    + \hat{v}_0(\xi) \frac{\sin(2\pi c \abs{\xi} t)}{2\pi c \abs{\xi}}.
\]
\end{theoreme}

\begin{proof}
La transform\'ee de Fourier en $x$ donne~:
$\hat{u}_{tt} = -(2\pi c)^2 \abs{\xi}^2 \hat{u}$. Pour chaque $\xi$ fix\'e,
c'est l'\'equation d'un oscillateur harmonique de pulsation $\omega = 2\pi c \abs{\xi}$.
La solution g\'en\'erale est $\hat{u}(\xi, t) = A(\xi)\cos(\omega t) + B(\xi)\sin(\omega t)$.
Les conditions initiales donnent $A = \hat{u}_0$ et $B = \hat{v}_0 / \omega$.
\end{proof}

\begin{exemple}[Formule de d'Alembert ($n = 1$)]
En dimension $1$, l'inversion donne~:
\[
    u(x,t) = \frac{u_0(x+ct) + u_0(x-ct)}{2} + \frac{1}{2c}\int_{x-ct}^{x+ct} v_0(s)\, ds.
\]
\end{exemple}

\section{Solutions fondamentales}

\begin{definition}[Solution fondamentale]
\index{solution fondamentale}
La \emph{solution fondamentale} (ou fonction de Green) d'un op\'erateur lin\'eaire
$L$ est la distribution $E$ telle que $LE = \delta$.
\end{definition}

\begin{exemple}[Laplacien en dimension $n \geq 3$]
La solution fondamentale de $-\Delta$ dans $\R^n$ ($n \geq 3$) est~:
\[
    E(x) = \frac{1}{n(n-2)\omega_n} \frac{1}{\abs{x}^{n-2}},
\]
o\`u $\omega_n = \frac{2\pi^{n/2}}{\Gamma(n/2)}$ est l'aire de $S^{n-1}$.
\end{exemple}

\begin{formules}[title=\textbf{Solutions fondamentales des EDP classiques}]
\begin{align*}
    \text{Chaleur~:} \quad & G_t(x) = \frac{1}{(4\pi \kappa t)^{n/2}}
        e^{-\abs{x}^2/(4\kappa t)}, \quad t > 0 \\[6pt]
    \text{Laplacien ($n=3$)~:} \quad & E(x) = \frac{1}{4\pi \abs{x}} \\[6pt]
    \text{Laplacien ($n=2$)~:} \quad & E(x) = -\frac{1}{2\pi}\ln\abs{x} \\[6pt]
    \text{Helmholtz ($n=3$)~:} \quad & E(x) = \frac{e^{ik\abs{x}}}{4\pi \abs{x}}
\end{align*}
\end{formules}

\section{Exercices}

\begin{exercice}
Calculer la transform\'ee de Fourier de $f(x) = e^{-a\abs{x}}$ pour $a > 0$
et $x \in \R$.
\end{exercice}

\begin{exercice}
Montrer la propri\'et\'e de convolution~: si $f, g \in L^1(\R^n)$, alors
$\widehat{f * g} = \hat{f} \cdot \hat{g}$.
\end{exercice}

\begin{exercice}
V\'erifier que le noyau de la chaleur $G_t$ satisfait~:
(a) $\int_{\R^n} G_t(x)\, dx = 1$ pour tout $t > 0$~;
(b) $(\partial_t - \kappa \Delta)G_t = 0$ pour $t > 0$~;
(c) $G_t \to \delta$ au sens des distributions quand $t \to 0^+$.
\end{exercice}

\begin{exercice}
Utiliser la transform\'ee de Fourier pour r\'esoudre l'\'equation de Schr\"odinger
libre~: $i\hbar \psi_t = -\frac{\hbar^2}{2m}\Delta \psi$, $\psi(x,0) = \psi_0(x)$.
\end{exercice}

\begin{exercice}
Soit $f \in \mathcal{S}(\R^n)$. Montrer que si $\hat{f}$ est \`a support compact,
alors $f$ se prolonge en une fonction enti\`ere sur $\C^n$ (th\'eor\`eme de
Paley-Wiener).
\end{exercice}

\begin{exercice}
R\'esoudre l'\'equation de la chaleur sur $\R$ avec donn\'ee initiale
$u_0(x) = \ind_{[0,1]}(x)$ (fonction indicatrice). Exprimer la solution
\`a l'aide de la fonction erreur $\operatorname{erf}$.
\end{exercice}

\begin{exercice}
Montrer que la solution fondamentale du laplacien dans $\R^n$ ($n \geq 3$) v\'erifie
$-\Delta E = \delta$ au sens des distributions, en calculant $\Delta E$ pour $x \neq 0$
et en utilisant le th\'eor\`eme de la divergence.
\end{exercice}
